\documentclass[12pt]{article}
\usepackage[a4paper,margin=1in]{geometry}
\usepackage[utf8]{inputenc}
\usepackage[T1]{fontenc}
\usepackage[hidelinks]{hyperref}
\pagestyle{empty}
\begin{document}
\begin{center}
{\Large\bfseries InkCloser: Line-Art Completion for Paint-Bucket Colorization\par}
\vspace{1em}
Ruben Tous\\
Computer Architecture Department, Universitat Polit\`ecnica de Catalunya, Barcelona, Spain\\
\texttt{ruben.tous@upc.edu}
\end{center}
\vspace{1.5em}
\noindent\textbf{Abstract}\par\medskip
\noindent Flood-fill (paint-bucket) coloring of line art requires the line drawing to form closed
regions. In practice, line art produced by drawing, scanning, or automatic line extraction
often contains small gaps, and a single gap lets the fill leak across the intended
boundary. We address the problem of restoring such gaps so that a
standard flood-fill tool can be applied without manual retouching. Unlike conventional
image inpainting, the location of the missing content is not available at inference
time. We present InkCloser, a U-Net with gated convolutions, a dilated
context stack, and transformer self-attention at two resolutions. We train it on pairs
built by synthetically removing segments of ink from a collection of clean, closed line-art
images, with a loss that combines pixel-wise cross-entropy, a global Dice term, and
a Dice term restricted to the neighborhood of the removed segments.
We also introduce evaluation metrics that measure gap closing and paint-bucket
behavior directly, rather than relying on pixel accuracy. On a held-out test set, the model closes 70.5\% of the gaps, raises the fraction of
regions that a single bucket click fills correctly (bucket accuracy) from 64.1\% to 94.7\%, reduces
the number of fill leaks per image from 15.5 to 1.5, and makes 45.0\% of the drawings
entirely leak-free. Code, pretrained weights and data are available at \url{https://github.com/rtous/inkcloser}.
\end{document}
